% Einheit 2 von 3 – Kegel (bank/pyramide-kegel-kugel/e2.jsonl, zusammenbau v0.1)
%% TODO Zweigzeile Teil 1: was hier gelernt wird (Satz in Schülersprache aus dem Einheitstitel) – Einheit 2
\einheitenkopf[e2]{Einheit 2 von 3 · Kegel}
\zweigzeile{ab Kl. 8, je nach Buch bis Kl. 10 · P10}
\setcounter{aufgabe}{17}

%% TODO Titel: Ich-kann-Satz fehlt in der Bank (Platzhalter: Kettenname) – Nr. 18
%% TODO Anweisung: ein Satz über der ersten Teilaufgabe fehlt in der Bank – Nr. 18
\begin{aufgabe}{Radius oder Durchmesser?}
\begin{teile}
\teil Eine Eistüte ist oben $6$ cm breit. Radius oder Durchmesser? Gib die andere Größe an. \leerfeld
\end{teile}
\end{aufgabe}

%% TODO Titel: Ich-kann-Satz fehlt in der Bank (Platzhalter: Kettenname) – Nr. 19
%% TODO Anweisung: ein Satz über der ersten Teilaufgabe fehlt in der Bank – Nr. 19
\begin{aufgabe}{Kegel}
\begin{teile}
\teil Die Strecke mit ? am Kegel – Höhe, Mantellinie oder Radius? Kreuze an.\\ \kreuz{Höhe}\\ \kreuz{Mantellinie}\\ \kreuz{Radius}
\end{teile}
%% TODO Darstellung für schwach fehlt in der Bank (pyramide-kegel-kugel-e2-k2-s1-v1; Abschnitt „Für schwache Schüler“)
\swz{Kegel: Radius $3$ cm, Höhe $7$ cm – Volumen? Runde auf eine Stelle.}{}
%% TODO Darstellung für schwach fehlt in der Bank (pyramide-kegel-kugel-e2-k2-s1-v2; Abschnitt „Für schwache Schüler“)
\swz{Kegel: Radius $6$ m, Höhe $4$ m – Volumen? Runde auf eine Stelle.}{}
%% TODO Darstellung für schwach fehlt in der Bank (pyramide-kegel-kugel-e2-k2-s1-v3; Abschnitt „Für schwache Schüler“)
\swz{Kegel: Radius $4$ cm, Höhe $6$ cm – Volumen? Runde auf eine Stelle.}{}
%% TODO Darstellung für schwach fehlt in der Bank (pyramide-kegel-kugel-e2-k2-s1-v4; Abschnitt „Für schwache Schüler“)
\swz{Kegel: Radius $2$ dm, Höhe $9$ dm – Volumen? Runde auf eine Stelle.}{}
%% TODO Darstellung für schwach fehlt in der Bank (pyramide-kegel-kugel-e2-k2-s1-v5; Abschnitt „Für schwache Schüler“)
\swz{Kegel: Radius $7$ cm, Höhe $15$ cm – Volumen? Runde auf eine Stelle.}{}
\swfrage{Was bleibt gleich, was ändert sich?}
%% TODO Darstellung für schwach fehlt in der Bank (pyramide-kegel-kugel-e2-k2-s2-v1; Abschnitt „Für schwache Schüler“)
\swz{Kegel: Durchmesser $8$ cm, Höhe $9$ cm – Volumen? Runde auf eine Stelle.}{}
%% TODO Darstellung für schwach fehlt in der Bank (pyramide-kegel-kugel-e2-k2-s3-v1; Abschnitt „Für schwache Schüler“)
\swz{Kegel: Radius $2{,}5$ cm, Höhe $6{,}4$ cm – Volumen? Runde auf eine Stelle.}{}
%% TODO Darstellung für schwach fehlt in der Bank (pyramide-kegel-kugel-e2-k2-s4-v1; Abschnitt „Für schwache Schüler“)
\swz{Ein Trichter ist innen ein Kegel mit dem Radius $6$ cm und der Höhe $11$ cm. Wie viel Milliliter fasst er?}{}
%% TODO Darstellung für schwach fehlt in der Bank (pyramide-kegel-kugel-e2-k2-s5-v1; Abschnitt „Für schwache Schüler“)
\swz{Kegel: Radius $16$ cm, Höhe $30$ cm – Mantellinie?}{}
%% TODO Darstellung für schwach fehlt in der Bank (pyramide-kegel-kugel-e2-k2-s6-v1; Abschnitt „Für schwache Schüler“)
\swz{Kegel: Mantellinie $29$ m, Radius $20$ m – Höhe?}{}
%% TODO Darstellung für schwach fehlt in der Bank (pyramide-kegel-kugel-e2-k2-s7-v1; Abschnitt „Für schwache Schüler“)
\swz{Kegel: Radius $4$ cm, Mantellinie $11$ cm – Mantel? Runde auf eine Stelle.}{}
%% TODO Darstellung für schwach fehlt in der Bank (pyramide-kegel-kugel-e2-k2-s8-v1; Abschnitt „Für schwache Schüler“)
\swz{Kegel: Radius $3$ cm, Mantellinie $7$ cm – Oberfläche? Runde auf eine Stelle.}{}
\swa{Skizziere das Schrägbild eines Kegels mit dem Radius $2$ cm und der Höhe $5$ cm. Zeichne die Höhe gestrichelt ein und beschrifte Radius und Höhe.}{\rechenplatz{5}}
%% TODO Darstellung für schwach fehlt in der Bank (pyramide-kegel-kugel-e2-k2-s10-v1; Abschnitt „Für schwache Schüler“)
\swz{Ein Kegel hat den Radius $3$ cm und die Mantellinie $8$ cm. Sein Netz besteht aus einem Kreis und einem Kreisausschnitt. Welchen Radius hat der Kreis, welchen der Kreisausschnitt?}{}
%% TODO Darstellung für schwach fehlt in der Bank (pyramide-kegel-kugel-e2-k2-s11-v1; Abschnitt „Für schwache Schüler“)
\swz{Kegel: Volumen $150$ cm³, Radius $4$ cm – Höhe? Runde auf eine Stelle.}{}
%% TODO Darstellung für schwach fehlt in der Bank (pyramide-kegel-kugel-e2-k2-s12-v1; Abschnitt „Für schwache Schüler“)
\swz{Ein Turm besteht aus einem $18$ m hohen Zylinder mit einem Kegeldach: Radius $3{,}5$ m, Mantellinie $6{,}2$ m. Das Dach bekommt Schiefer für $45$ € je m². Was kostet das, und wie hoch ist das Bauwerk insgesamt \hfill (P10 2026 FOR)?}{}
%% TODO Darstellung für schwach fehlt in der Bank (pyramide-kegel-kugel-e2-k2-s13-v1; Abschnitt „Für schwache Schüler“)
\swz[3]{Ein Kegel und ein Zylinder haben beide den Radius $3$ cm und die Höhe $5$ cm. Berechne beide Volumen. Wie oft passt der Kegel in den Zylinder?}{}
\swz[3]{Ein Kegel und ein Zylinder sind gleich hoch. Mia behauptet: „Ist der Radius des Kegels sechsmal so groß wie der Radius des Zylinders, hat der Kegel das doppelte Volumen.“ Prüfe die Aussage mit einer Rechnung \hfill (P10 2026 FOR).}{}
\end{aufgabe}

%% TODO Titel: Ich-kann-Satz fehlt in der Bank (Platzhalter: Kettenname) – Nr. 20
%% TODO Anweisung: ein Satz über der ersten Teilaufgabe fehlt in der Bank – Nr. 20
\begin{aufgabe}{Kegel – Fehler finden, Begründen, Darstellungswechsel, Anwendung}
\swz[3]{Kegel: Radius $4$ cm, Höhe $15$ cm. Paul rechnet das Volumen: \rechnung{V &= \pi \cdot 4^2 \cdot 15 \\ V &\approx 754{,}0 \text{ cm}^3} Wo ist der Fehler?}{}
\swfrage{Ein Kegel und ein Zylinder haben gleichen Radius und gleiche Höhe. Begründe, warum man den Kegel dreimal füllen muss, bis der Zylinder voll ist.}
\swa{Eine Waffeltüte ist oben $5$ cm breit und $11$ cm tief. Skizziere sie als Kegel mit der Spitze unten und beschrifte Radius und Höhe.}{\rechenplatz[halb]{4}}
\swz[3]{Auf dem Bauhof liegt ein kegelförmiger Haufen Streusplitt: Durchmesser $6$ m, Höhe $2$ m. Ein Lkw kann $15$ m³ laden. Reicht eine Fuhre, um den Haufen abzufahren?}{}
\end{aufgabe}

\uebersichtskasten{Volumen: ein Drittel von Grundkreis mal Höhe – der Zylinder mit gleichem Radius und gleicher Höhe fasst dreimal so viel. \\ r = 5 cm, h = 12 cm: V = π · 5² · 12 : 3 = 100 · π ≈ 314,2 cm³ \quad Zylinder: π · 5² · 12 ≈ 942,5 cm³ \\ Mantellinie: rechtwinkliges Dreieck aus Höhe und Radius, die Mantellinie ist die Hypotenuse. \quad s = √(12² + 5²) = 13 cm \quad umgekehrt h = √(s² − r²) \\ Mantel: π mal Radius mal Mantellinie – kein Rechteck wie beim Zylinder, sondern ein Kreisausschnitt. \quad M = π · 5 · 13 ≈ 204,2 cm² \quad Oberfläche: O = π · 5² + M ≈ 78,5 + 204,2 = 282,7 cm² \\ Rückwärts: h = 3 · V : (π · r²).}
